\paragraph{The alignment beats linear CCA; the sparse code is for interpretation.} Table~\ref{tab:main}. Our crosscoder's linear alignment head retrieves the paired variant at Recall@1 0.307$\pm$0.014 (5 seeds, residue-disjoint $n{=}\posNtest{}$), significantly above linear CCA (0.216; paired bootstrap $\Delta{=}0.104$, $p<0.001$) and matching a pure deep-CCA baseline (0.311). The \emph{sparse} shared code alone does not retrieve (0.003, chance) --- the cross-modal signal lives in the learned linear alignment, and the sparse codes serve interpretability (below). The same ordering holds for DMS (shared 0.442$\pm$0.003 $>$ CCA 0.337, DNA 0.375; ESM-2 0.514 and concat 0.514 remain stronger). Reconstruction FVE is 0.72 (DNA)/0.70 (protein).
\begin{table}[t]\centering\small
\caption{Cross-modal retrieval and DMS prediction (residue-disjoint, $n{=}\posNtest{}$; retrieval/DMS from 5 seeds). A learned alignment (deep-CCA or our crosscoder's alignment head) beats linear CCA at both; the crosscoder's \emph{sparse} shared code does not retrieve (it is the interpretable dictionary, not the alignment). ESM-2 alone / concat are the strongest DMS predictors.}\label{tab:main}
\begin{tabular}{lccc}
\toprule
Method & R@1 & R@10 & DMS $\rho$ \\
\midrule
linear CCA & 0.216 & 0.645 & 0.337 \\
deep CCA (align only) & 0.311 & 0.689 & 0.464 \\
\;crosscoder: sparse shared code & 0.003 & 0.027 & 0.272 \\
\;crosscoder: alignment head & 0.307 & 0.689 & 0.442 \\
ESM-2 $\Delta$ (probe) & -- & -- & 0.514 \\
concat / PLS & -- & -- & 0.514 \\
\bottomrule
\end{tabular}
\end{table}

\paragraph{Domain-disjoint generalization.} Training and evaluating only across the two disjoint domains (train BRCT, test RING; $n{=}\domNtest{}$), the alignment still beats CCA (Recall@1 0.106 vs.\ 0.081; DMS 0.303 vs.\ 0.279), though with substantial degradation (ESM-2 alone reaches 0.444). The method partially transfers to an unseen structural domain rather than memorizing residues.

\paragraph{Cross-gene generalization (multi-gene ClinVar).} To test the single-gene concern directly, we assemble \mgNgenes{} genes of ClinVar missense variants (balanced pathogenic/benign, GRCh38 coordinates, protein references validated) and train on \mgNtrainGenes{} genes, evaluating on \mgNtestGenes{} \emph{held-out} genes ($n{=}\mgNtest{}$). Cross-modal retrieval within each unseen gene reaches Recall@1 0.066 (Recall@10 0.340), above linear CCA (0.056) and comparable to deep-CCA (0.068); the alignment thus transfers to genes never seen in training. On \emph{gene-disjoint} ClinVar pathogenicity prediction, the shared code reaches AUROC 0.886, versus the Evo\,2 probe 0.864, the ESM-2 probe 0.855, and CCA 0.895---the shared cross-modal representation carries transferable pathogenicity signal across genes.

\paragraph{Biological structure of the sparse codes.} On held-out variants the shared sparse code predicts loss-of-function (AUROC 0.830, 95\% CI [0.776,0.877], $n{=}\lofN{}$) and domain (0.839, $n{=}\ringN{}$). Modality-private codes are complementary: the protein-private code is the best domain predictor (0.853) and matches shared on LOF (0.842), consistent with ESM-2 encoding structural context. The shared code also exceeds classical single-variant VEP predictors (SIFT 0.305, CADD 0.329, phyloP 0.336; $|\rho|$ vs.\ DMS).

\paragraph{A shared functional axis.} The direction of maximal DMS correlation, fit independently in each model's activation space, maps to nearly the same point in the alignment space (cosine 0.905), far above random direction pairs (0.306, 95th pctile $|{\cos}|=0.513$). This is property-specific: the two \emph{domain} directions also align (0.927), while a functional-vs-domain pair sits at the random level (0.282). Same biological property $\Rightarrow$ same shared axis across modalities; different properties $\Rightarrow$ different axes.

\paragraph{Causal probing (a negative result).} We test whether this alignment yields a causal handle by injecting the shared functional direction at the variant position and re-scoring ($n{=}\causalN{}$ held-out variants). Single-position activation edits do \emph{not} give reliable causal control in either model: the score shift is uncorrelated with the variant's true effect and indistinguishable from a matched-norm random direction (ESM-2 Spearman 0.131, $p{=}0.19$; Evo\,2 0.009, $p{=}0.93$). Representational alignment does not imply an interchangeable causal handle---a caution for cross-model steering, and an open question of whether richer interventions would succeed.